\subsection{Smooth Functions on Manifolds}
\lecdate{Lec 7 - Jan 27 (supposedly) (Week 6)}

principles:
\begin{enumerate}
    \item charts and their inverses are smooth (by defn)
	\item compositions of smooth fns are smooth
\end{enumerate}
the [properties and behaviour we want] follow largely from these principles.

\begin{defn}
    a fn $f:M\gto\bR$ is \textbf{smooth at} $\bm{p}$ if for some chart
	$(U,\vphi)$ with $p\in U$, the map $f\circ\vphi^{-1}:\vphi(U)\gto\bR$
	is smooth (i.e. inf-diff'ble) at $\vphi(p)$.

	we say $f$ is \textbf{smooth} if it is smooth at every pt.
	equivalently, $f\circ\vphi^{-1}$ is smooth for every $\vphi$,
	or simply some $\vphi$ which cover $M$ (easiest to check).
\end{defn}
\textbf{remark}: it doesnt depend on what chart you pick, since
transition maps are smooth.
notably, we can replace ``some" w ``every" and nothing changes.

\begin{xmp}[source=Primary Source Material]
    \vspace{-0.2in}
	\begin{enumerate}[1)]
		\item if $f:\bR^{n+1}\gto\bR$ smooth, then
			$f\big\rvert_{S^{n}}:S^{n}\subseteq\bR^{n+1}\gto\bR$ smooth.

			indeed, take stereo proj. then $\vphi_{N}^{-1},\vphi_{S}^{-1}$ smooth,
			and the rest follows.
		\item $f:S^{2}\gto\bR$, given as $f(x,y,z)=\sqrt{1-z^{2}}=\sqrt{x^{2}+y^{2}}$.
			using the chart $\vphi:\set{z>0}\gto B(0,1)$ as $\vphi(x,y,z)=(x,y)$,
			want $(f\circ\vphi^{-1})(x,y)=\sqrt{x^{2}+y^{2}}$ smooth.
			but $f\circ\vphi^{-1}$ not smooth at $(0,0)$, so $f$ not smooth.
	\end{enumerate}
\end{xmp}
p sure he messed up the 2nd example, should have $f(x,y,z)=z$ so $(f\circ\vphi^{-1})(x,y)=\sqrt{1-(x^{2}+y^{2})}$.
tis what book says anyway.

\begin{defn}
    we define $C^{\infty}(M)=\set{\trm{smooth fn's }M\gto\bR}$.
\end{defn}

some facts:
\begin{itemize}
	\item if $f,g:M\gto\bR$ smooth, then $af+bg$ smooth where $a,b\in\bR$.
	\item products of smooth fn's are smooth
	\item the constant fn $1\in C^{\infty}(M)$ is smooth
\end{itemize}
in other words, $C^{\infty}(M)$ is an $\bR$-algebra (lmao?).

\begin{defn}
    cts fn's on $M$ are defn'd similarly:
	$f:M\gto\bR$ cts if $f\circ\vphi^{-1}$ cts for all (equivalently, enough) charts $\vphi$.
	in particular, smooth fn's are cts.
\end{defn} \

\begin{prop}
    $f:M\gto\bR$ is cts iff $f^{-1}(J)$ open for all open $J\subseteq\bR$.
\end{prop} \

\begin{pf}[source=Primary Source Material]
	by the same fact for $U\subseteq\bR^{n}$, we have:
	\begin{equation*}
		(f\circ\vphi_{\alpha}^{-1})^{-1}(a,b) \subseteq \vphi_{\alpha}(U_{\alpha})
	\end{equation*}
	open for every chart $\vphi_{\alpha}$.
	since $V\subseteq M$ open iff it's a union of $\vphi_{\alpha}^{-1}(V_{\alpha})$
	for open $V_{\alpha}\subseteq\vphi_{\alpha}(U_{\alpha})$, then:
	\begin{equation*}
	    f^{-1}(a,b)=\bigcup_{\alpha}
		\vphi_{\alpha}^{-1}((f\circ\vphi_{\alpha}^{-1})(a,b))
	\end{equation*}
	is open. [huh?]
\end{pf} \

\begin{defn}
    the \textbf{support} of $f:M\gto\bR$ is the set:
	\begin{equation*}
	    \supp(f) \ := \ \bar{f^{-1}(\bR\sm\set{0})}
	\end{equation*}
	equivalently, if $p\notin\supp(f)$, then $f$ is 0 in a nbhd of $p$.
\end{defn} \

\begin{lm}
    sps $U\subseteq M$ open, $g:U\gto\bR$ smooth, $\supp(g)\subseteq U$
	closed in $M$.
	then $g$ extends to a smooth $f:M\gto\bR$ where
	$f\big\rvert_{U}=g$ and $f(M\sm U)=0$.
\end{lm} \

\begin{pf}[source=Primary Source Material]
    let $V=M\sm\supp(g)$ open in $M$.
	then $f$ is smooth at every pt $p\in U$ and $V$, since we can choose
	chart containing $p$ w domain in $U$ or $V$ resp., and $f\equiv0$ in the case of $V$.
	since $M=U\cup V$, we're done.
\end{pf}

im gonna assume we did 3.2, which talks abt $M$ being T2 iff pts can be separated by smooth fn's.
pf at an intuitive lvl is basically that $M$ is ``locally metrizable", but
details of fwd dir is to construct $f$ using bump fn's.

replacing w cts also loses the if dir;
presumably can happen on mflds w weird topologies.

\lecdate{Lec 8 - Jan 29 (Week 4)}
whered my tuesday notes go.

recall that a map $f:M\gto N$ btwn 2 mflds is smooth at a pt $p\in M$
if $\ex$ charts $(U,\vphi)$ for $M$ and $(V,\psi)$ for $N$ s.t.
$\psi\circ f\circ\vphi^{-1}:\vphi(U\cap f^{-1}(V))\gto\psi(V)$ is smooth at $\vphi(p)$.
equivalently, just require that $f(U)\subseteq V$.
its smooth if its smooth at all pts, and a diffeo if it and its inv are smooth.

intuitively: this means $f$ is smooth as a real fn after placing coords on $M$ and $N$.

last time, showed $\bR P^{1}\gto S^{1}$ given by $(x:1)\mto\vphi_{N}^{-1}(x)$
and $(1:0)\mto(N)$ is diffeo.

\begin{xmp}[source=Primary Source Material]
    \vspace{-0.2in}
	\begin{enumerate}[(1)]
		\item $\bR^{n+1}\sm\set{0}\gto\bR P^{n}$ via proj $\pi$, w
			$(x^{0},\dots,x^{n}) \mto (x^{0}:\cdots:x^{n})$. smooth since:
			\begin{gather*}
			    p=(x^{0},\dots,x^{n}) \xrightarrow{\pi} (x^{0}:\cdots:x^{n}) \xrightarrow{\vphi_{i}} (1/x^{i})(x^{0},\dots,x^{i-1},x^{i+1},\dots,x^{n}) \\
			    (x^{0},\dots,x^{n}) \xrightarrow{\trm{id}} (x^{0},\dots,x^{n}) \mto (1/x^{i})(x^{0},\dots,x^{i-1},x^{i+1},\dots,x^{n})
			\end{gather*}
		\item $\bC^{n+1}\sm\set{0}\xrightarrow{\pi}\bC P^{n}$ w $(z^{0},\dots,z^{n})\mto(z^{0}:\cdots:z^{n})$.
			same pf.
		\item $\bC P^{1}\gto S^{2}$, w $(z=x+iy:1)\mto\vphi_{N}^{-1}(x,y)$
			and $(1:0)\mto N$. smooth:
			\begin{equation*}
				p=(z:1) \qquad \vphi_{N}(f(p)) = (x,y) = \trm{id}(\vphi_{1}(p))
			\end{equation*}
			for $(1:0)$:
			\begin{gather*}
				\bC P^{1}\sto\bC \qquad \vphi_{0}(z:1)=\vphi(1:1/z)=1/z=\bar{z}/\abs{z}^{2}
				\qquad \vphi_{0}(1:0)=0 \\
				S^{2}\sto\bR^{2} \qquad \vphi_{S}\circ\vphi_{N}^{-1}(x,y)=\frac{(x,y)}{x^{2}+y^{2}}=\frac{z}{\abs{z}^{2}}
				\qquad \vphi_{S}(N)=0
			\end{gather*}
		\item $f:\bR\gto\bR$ given by $f(x)=x^{3}$ is smooth + bijective,
			but $f^{-1}(x)=x^{1/3}$ not diff'ble at $x=0$.
	\end{enumerate}
\end{xmp}
rmks: check that composition of smooth is smooth.
also, diffeo is the ``right" notion of equiv btwn smooth mflds; that is,
properties defn'd from smooth structure are preserved under diffeos.

Q: are homeo mflds always diffeo?

A: surprisingly, yes in dim 1,2,3, but not higher dims!
there are 28 smooth structures on $S^{7}$, known as ``exotic spheres" (LMFAO),
found by Milnor in the 50s.

in 80s, found $\infty$ many spaces homeo but not diffeo to $\bR^{4}$
(freedman, donaldson, etc.).

consider the following composition:
\begin{equation*}
    S^{3}\ito\bR^{4}\sm\set{0}\sto\bC P^{1}\simeq S^{2}
\end{equation*}
given by inclusion then projection.
the composition map $h$ is known as the \textbf{Hopf fibration}.
what can we say abt it?
\begin{enumerate}[(1)]
	\item its smooth
	\item given $p:=(z:w)\in S^{2}$, what does $h^{-1}(p)$ look like?
		can assume that $\abs{z}^{2}+\abs{w}^{2}=1$. answer:
		$h^{-1}(z:w)=\set{(e^{i\theta}z,e^{i\theta}w)}\subseteq S^{3}$ is a circle.
\end{enumerate}
which circles?

recall $S^{3}=\bR^{3}\cup\set{\infty}$, viewed via stereo proj.
in this case, $(z_{1},z_{2},w_{1},w_{2})\mto(1-w_{2})^{-1}(z_{1},z_{2},w_{1})$.
we have:
\begin{gather*}
    h^{-1}(0:1)\gto(0,0,t) \qquad t\in\bR \\
	h^{-1}(1:0)\gto(\cos\theta,\sin\theta,0) \\
	h^{-1}(z:r\in\bR)\gto
	\left(\frac{e^{i\theta}z}{1-r\sin\theta},\frac{r\cos\theta}{1-r\sin\theta}\right)
\end{gather*}
he drew a cool picture that i cant draw rn :(



